\section{功能性需求}

\subsection{执行者分析}

执行者是位于系统边界之外、通过系统获得业务价值或向系统提供服务的人或外部系统。图~\ref{fig:actor-analysis} 说明执行者之间的泛化与协作关系；同一个自然人在通过审核后可以拥有多个业务角色，但每次请求仍按当前角色及资源归属授权。

\srsfigure{figures/srs/09-actor-analysis.pdf}{系统执行者分析图}{UML 用例图}{fig:actor-analysis}

\begin{longtable}{|L{0.16\linewidth}|L{0.26\linewidth}|L{0.48\linewidth}|}
\caption{系统执行者说明}\label{tab:actors}\\ \hline
\textbf{执行者} & \textbf{执行者关系} & \textbf{与系统的主要交互} \\ \hline
游客 & 独立执行者 & 浏览公开店铺与商品、注册账号、使用公开规则问答。 \\ \hline
已登录用户 & 顾客、商家、骑手和管理员的通用执行者 & 维护本人资料、查看通知、切换已获授权的角色。 \\ \hline
顾客 & 已登录用户的特化角色 & 管理地址和购物车，创建、支付、取消、跟踪和评价本人订单，使用个人资产及智能点餐。 \\ \hline
商家 & 已登录用户的特化角色 & 申请经营资格，维护本人店铺商品，处理本店订单、回复评价并查看经营数据。 \\ \hline
骑手 & 已登录用户的特化角色 & 申请配送资格，上下线、接单、到店、取餐、送达和上报异常。 \\ \hline
管理员 & 已登录用户的特化角色 & 完成准入审核、账号与内容治理、配送异常处理和平台统计。 \\ \hline
外部服务 & AI、地图和对象存储等支撑执行者 & 通过受控接口提供生成、导航和媒体存储能力，不直接改变核心业务状态。 \\ \hline
\end{longtable}

\subsection{总用例}

\srsfigure{figures/srs/10-system-use-case.pdf}{系统总用例图}{UML 用例图}{fig:system-use-case}

总用例图展示游客、四类登录角色和外部服务与一级用例的关系。角色用例在后续分图和用例表中展开，需求编号同时用于接口、数据库和测试材料。
